\documentclass[11pt]{article}
\usepackage[margin=1in]{geometry}
\usepackage{amsmath,amssymb,amsthm,booktabs,array}
\usepackage[T1]{fontenc}
\usepackage{lmodern,microtype}
\DeclareMathSizes{10}{10}{7}{7}
\DeclareMathSizes{10.95}{10.95}{8}{8}
\usepackage[colorlinks=true,linkcolor=blue,urlcolor=blue,citecolor=blue]{hyperref}
\hypersetup{pdftitle={Sharp Newton-edge bounds and cancellation families with one nonzero quadratic norm},
pdfauthor={Ying Xie},pdfsubject={Newton polygons over valued fields}}
\newtheorem{theorem}{Theorem}[section]
\newtheorem{lemma}[theorem]{Lemma}
\newtheorem{corollary}[theorem]{Corollary}
\theoremstyle{remark}\newtheorem{remark}[theorem]{Remark}
\newcommand{\val}{\nu}
\newcommand{\qf}{\mathfrak q}
\newcommand{\Q}{\mathbb Q}
\newcommand{\Z}{\mathbb Z}
\newcommand{\R}{\mathbb R}
\newcommand{\supp}{\operatorname{supp}}
\title{Sharp Newton-edge bounds and cancellation families\newline
with one nonzero quadratic norm}
\author{Ying Xie\\Kennesaw State University\\
\href{mailto:yxie2@kennesaw.edu}{yxie2@kennesaw.edu}}
\date{}
\begin{document}
\maketitle
\begin{abstract}
Let $g=(P,Q,R)=\sum_{i=0}^{5}c_i x^{A_i}$ have six occupied vector
coefficients, exactly one of which has nonzero norm for
$\qf(P,Q,R)=PQ+R^2$. Over every field of characteristic different from
two with a real-valued nonarchimedean valuation, we prove that
$\qf(g)$ has at most eleven strict lower Newton edges and at most ten
edges of width one. Both maxima are attained simultaneously over every
such nontrivially valued field. At every prime, integer examples have
coordinate sparsity fourteen and exactly ten simple nonzero local roots
of distinct valuations. The dyadic example has 228-bit input coefficients.
The upper bound is computer assisted: a generic null-cone perturbation
and a rank-three initial-matrix argument reduce the problem to 2,359
marked support patterns. Exact disjunction trees and integer linear
identities exclude every pattern, without an exponent or precision
cutoff. Six patterns require convexity inequalities forced by equal
barycenters of exponent sums. The supplied verifier uses only the Python
standard library. We also prove the sharp five-support bound of eight
edges and construct, at every prime and every support size $n\ge8$,
one-norm examples with $2n+1$ edges, of which $2n$ have width one.
The envelope lemmas and the family construction are proved in full.
\end{abstract}

\section{The extremal problem}

Let $K$ be a field of characteristic different from two and let
$\val:K\to\R\cup\{+\infty\}$ be a nonarchimedean valuation. For
$f=\sum_e f_e x^e\ne0$, write $V(f)$ for the number of strict edges on
the lower convex hull of the finite points $(e,\val(f_e))$, and $U(f)$
for the number of those edges having horizontal width one. Collinear
coefficient points do not create additional edges. Laurent exponents are
allowed. Put
\[
 \qf(p,q,r)=pq+r^2,\qquad
 g=(P,Q,R)=\sum_{i=0}^{n-1}c_i x^{A_i},\quad A_0<\cdots<A_{n-1}.
\]
All $c_i$ are nonzero. Let $k=|\{i:\qf(c_i)\ne0\}|$, and let $S$ be
the sum of the three scalar support sizes.

\begin{theorem}\label{thm:main}
If $n=6$, $k=1$, and $\qf(g)\ne0$, then
\[
 V(\qf(g))\le11,\qquad U(\qf(g))\le10.
\]
Over every $K$ as above with nontrivial valuation, a formula attains
$(V,U)=(11,10)$. No bound on the input exponents or coefficient
valuations is imposed. No completeness or residue-field assumption is needed.
\end{theorem}

\begin{corollary}\label{cor:prime}
For every prime $p$ there are integer polynomials $P,Q,R$ with
$n=6$, $k=1$, $S=14$, and $(V_p,U_p)=(11,10)$. Their sole nonzero
coefficient norm is an ordinary integer square. The output has exactly
ten simple nonzero $\Q_p$ roots, of distinct valuations, and a separate
double zero at the origin.
\end{corollary}

The small-support results fit into a wider cancellation picture.
Section~\ref{sec:five} proves that five coefficient vectors with one
nonzero norm have at most eight edges, and gives every-prime examples
with eight unit edges. Together with Theorem~\ref{thm:main}, this makes
six the smallest support size violating $V\le2n-2$ over $\Q_p$.
Section~\ref{sec:family} constructs integer examples with
$V_p=2n+1$, $U_p=2n$, and $S=3n-4$ for every $n\ge8$ and every prime.
Their input bit lengths are $O(n^2\log p)$. Thus the small-support
extrema coexist with larger families exceeding $2n$ edges.
All envelope arguments and constructions used for these statements
are included below.

Sparse Newton-polygon bounds are motivated in part by algebraic
complexity. Koiran, Portier, Tavenas, and Thomass\'e study a
$\tau$-conjecture for ordinary bivariate Newton polygons and show that
a suitable weak version implies an arithmetic-circuit lower bound for
the permanent \cite{KPTT}. Their polygon and representation model
are different from the valued univariate, fixed ternary quadratic
problem here. We use their work only as context; we do not infer its
complexity conclusion from a six-support bound.

\section{Reduction to finite pair and coefficient heights}

Write $c_i=(P_i,Q_i,R_i)$ and set
\[
 d_i=\qf(c_i),\qquad
 C_{ij}=P_iQ_j+P_jQ_i+2R_iR_j \quad(i<j).
\]
The matrix $M$ with diagonal $2d_i$ and off-diagonal $C_{ij}$ has rank
at most three, and
\[
 f=\qf(g)=\sum_i d_i x^{2A_i}+\sum_{i<j}C_{ij}x^{A_i+A_j}.
\]
Let $a$ be the unique index with $d_a\ne0$.

\begin{lemma}[Generic null-cone perturbation]\label{lem:perturb}
After passing to a real-valued nonarchimedean field extension, the five
null vectors can be perturbed so that every $C_{ij}$ and every possible
output coefficient is nonzero. The support and the number of nonzero
coefficient norms remain unchanged, and every original strict Newton
vertex can be preserved. Every original unit edge can also be preserved.
\end{lemma}
\begin{proof}
Every nonzero null vector has the form
$t_i(-u_i^2,v_i^2,u_iv_i)$ with $t_i\ne0$. If $Q_i\ne0$, take
$t_i=Q_i$, $v_i=1$, and $u_i=R_i/Q_i$. If $Q_i=0$, nullity gives
$R_i=0$ and $P_i\ne0$; use $t_i=-P_i$, $u_i=1$, $v_i=0$.

Replace $u_i,v_i$ by $u_i+Z_i,v_i+W_i$, where the ten variables are
algebraically independent and have a common sufficiently large positive
Gauss valuation. Keep $c_a$ fixed. The perturbed vectors are nonzero
and null. Between two null vectors the cross coefficient is
\[
 -t_it_j(u_iv_j-u_jv_i)^2,
\]
which is a nonzero polynomial in the independent parameters. With
$c_a$ it is the nonzero quadratic form
$t_i(-Q_a u_i^2+P_a v_i^2+2R_a u_iv_i)$.

Pairs at a common exponent have disjoint endpoints. A null-null pair
has a mixed monomial involving its own two endpoints, absent from all
other contributions in that exponent group. A pair involving $a$ has
a nonconstant polynomial in its own null endpoint, likewise absent
elsewhere in the group. Thus no possible output coefficient vanishes
identically, even if a diagonal constant is present.

For each old strict vertex select a finite supporting slope exposing
it uniquely. There are finitely many selected slopes and possible
output exponents. Sufficiently large parameter valuations preserve
all old nonzero coefficient heights and place every newly nonzero
coefficient strictly above these selected supporting lines. All old
vertices therefore persist. A unit edge persists because no integer
exponent lies between its endpoints.
\end{proof}

The use of a field extension is legitimate for the universal upper
bound: the alleged counterexample would become another counterexample
over a field in the same class. The lemma does not assert that a
specialization can always be chosen inside the original field.

\section{The hidden primary pair}

For $t\in\R$ define the concave Gauss functions
\begin{align*}
 E(t)&=\min\{\val(d_i)+2A_it,\ \val(C_{ij})+(A_i+A_j)t\},\\
 D(t)&=\min_i\{\val(d_i)+2A_it\},\\
 F(t)&=\min_e\{\val(f_e)+et\}.
\end{align*}
Zero contributions are assigned infinite height. The following envelope
lemma holds for arbitrary $n$ and $k$; its full proof supplies the
preliminaries for both small-support bounds.

\begin{lemma}\label{lem:envelope}
For rank at most three, $E=\min(F,D)$. If $E=F$, then
$V(f)\le2n+k-4$ for $n\ge2$.
\end{lemma}
\begin{proof}
Always $E\le\min(F,D)$. If the inequality were strict on an interval,
choose $t$ there avoiding ties between distinct contribution exponents.
Over $K(T)$ with the Gauss valuation $\val(T)=t$, conjugate $M$ by
$\operatorname{diag}(T^{A_i})$ and divide by a minimum off-diagonal
entry times its monomial. All entries are integral, all diagonal
reductions vanish, and the nonzero off-diagonal reductions are precisely
the minimum pairs. There must be at least two such pairs, since the
output coefficient is raised. Equal-sum pairs are disjoint, so the
residue matrix contains two nonsingular two-by-two blocks and has rank
at least four. This contradicts the vanishing of every four-by-four
minor. Continuity gives the identity at exceptional parameters too.

Now assume $E=F$. At each strict output vertex, choose a contribution
of that height. Such a contribution exists by the ultrametric inequality.
At most $k$ selected vertices use diagonals. Two strictly nested chosen
off-diagonal pairs, on four ordered exponents $x_0<x_1<x_2<x_3$, are
impossible. Here is an explicit strictness argument. Write
$s=x_0+x_3$, $t=x_1+x_2$, and let $\phi$ be the piecewise-linear
Newton boundary, extended convexly with strict endpoint corners if needed.
The selected reverse entries have heights $\phi(s),\phi(t)$, and
both arguments are strict vertices. For sufficiently small $\eta>0$,
the function
$\psi(x)=\phi(x)-\eta(|x-s|+|x-t|)$ is still convex.
For any convex function, the reverse permutation minimizes the assignment
sum $\sum_i\psi(x_i+x_{\pi(i)})$: successively swapping increasing
permutation values cannot increase the sum, by the convex Monge inequality.

Every other permutation has an argument outside the interval with
endpoints $s,t$. Indeed, $\pi(0)\ne3$ gives an argument below both;
$\pi(3)\ne0$ gives one above both. The remaining nonreverse permutation
fixes the middle indices, and one of $2x_1,2x_2$ lies outside that
interval. Therefore its sum of $|x-s|+|x-t|$ is strictly larger than
the reverse sum. The reverse permutation uniquely minimizes the
$\phi$ assignment cost. All matrix entries are at least their boundary
heights, since $E=F$ (and $\val(2)\ge0$ for doubled diagonals), while
the reverse entries attain their heights. Thus the reverse determinant
term is uniquely minimum, contradicting rank at most three.

The selected pairs consequently form a nonnesting antichain. In
lexicographic order both indices are nondecreasing and their sum
strictly increases, through the range $1,\ldots,2n-3$ in zero-based
indexing. There are at most $2n-3$ pairs. Adding $k$ diagonal vertices
gives $V+1\le2n-3+k$.
\end{proof}

\begin{corollary}\label{cor:visible-diagonals}
For $n\ge2$, if every nonzero diagonal contribution lies on or above
the output Newton boundary, then $V(f)\le2n+k-4$. In particular,
$V(f)\le2n-4$ when all coefficient vectors are null. The bound
$V(f)\le2n+k-4$ also holds whenever the support contains no nontrivial
three-term arithmetic progression.
\end{corollary}
\begin{proof}
The diagonal hypothesis says $D\ge F$, so Lemma~\ref{lem:envelope}
gives $E=F$ and the bound. For null vectors $D=+\infty$.
If the support has no nontrivial three-term progression, no
off-diagonal pair contributes at $2A_i$; hence $f_{2A_i}=d_i$ and
every finite diagonal point belongs to the output coefficient cloud.
\end{proof}

\begin{lemma}\label{lem:primary}
For $n=6,k=1$, any formula with $V>9$ can be normalized to have a
unique primary pair $(p,q)$ satisfying
\begin{gather*}
 p<a<q,\qquad A_p+A_q=2A_a,\\
 w_{pq}=0,\qquad w_{ij}>0\ ((i,j)\ne(p,q)),\qquad H_e>0.
\end{gather*}
Here $w_{ij}$ are pair-entry heights and $H_e$ are output heights
after a common vertical translation and affine tilt. The normalized
exceptional diagonal contribution has height zero. Reflection permits
$a\in\{1,2\}$.
\end{lemma}
\begin{proof}
Lemma~\ref{lem:perturb} permits all pair and output coefficients to be
nonzero. If $E=F$, Lemma~\ref{lem:envelope} gives $V\le9$. Thus $E<F$
on an open interval. At a generic parameter in that interval, the
unique diagonal at exponent $2A_a$ is minimal, and at least one pair
at the same exponent must tie it. Such a pair avoids $a$ and straddles it.

There cannot be two minimum pairs. They would give two disjoint
nonsingular blocks in the residue matrix used above. With
$\epsilon=\val(2)>0$ the diagonal vanishes on reduction; with
$\epsilon=0$ it adds a separate nonzero diagonal block. Either case
contradicts rank at most three. Subtracting the minimum height and
tilting by the chosen parameter gives the assertions. The extreme
indices cannot be exceptional in a hidden primary, and reflection
reduces $a=3,4$ to $a=2,1$.
\end{proof}

\section{The sharp five-support bound}\label{sec:five}

\begin{theorem}\label{thm:five}
If $n=5$, $k=1$, and $\qf(g)\ne0$, then $V(\qf(g))\le8$ over
every field in the standing hypotheses. For every prime $p$, integer
examples attain $V_p=U_p=8$ with $S=11$ and a sole nonzero coefficient
norm that is an integer square.
\end{theorem}

\begin{proof}[Proof of the upper bound]
Suppose $n=5$, $k=1$, and, toward a contradiction, $V\geq9$.
Let $a$ be the exceptional exponent. There are only eleven possible
contributions: ten off-diagonal pairs and $d_a x^{2a}$.
Lemma~\ref{lem:envelope} implies $E\neq F$, and Lemma~\ref{lem:envelope}
then forces the diagonal to be hidden by cancellation.
Counting contributions shows that all ten off-diagonal terms are
nonzero with distinct exponent sums, precisely one pair
$p<a<q$ satisfies $p+q=2a$, and every one of the ten output points is
a strict vertex. No output coefficient vanishes.

Put $d=\val(d_a)$, $w_{ij}=\val(C_{ij})$, and let $h_{ij}$ be the
actual output height at $i+j$. In this section the letters $p,a,q,r,s$ denote actual support exponents, not their ordinal indices. Then
\begin{equation}\label{five:eq:onecollision}
 w_{pq}=d<h_{pq},\qquad w_{ij}=h_{ij}\quad(\{i,j\}\neq\{p,q\}).
\end{equation}
Strict convexity gives, whenever $i<j<k<l$,
\begin{equation}\label{five:eq:convex4}
 h_{ij}+h_{kl}>h_{ik}+h_{jl}>h_{il}+h_{jk},
\end{equation}
and, for $r\notin\{p,a,q\}$,
\begin{equation}\label{five:eq:midpoint}
 h_{pr}+h_{qr}>2h_{ar}.
\end{equation}
In \eqref{five:eq:convex4} the paired arguments have equal totals and
successively smaller spreads.

For four null vectors, the Gram determinant identity is
\begin{equation}\label{five:eq:gramnull}
 X^2+Y^2+Z^2-2XY-2XZ-2YZ=0,
\end{equation}
where $X,Y,Z$ are the three matching products of cross coefficients.
Their valuations cannot have a unique minimum.
For the exceptional vector $a$ and three null vectors $p,q,r$, the
identity has the additional term
\begin{equation}\label{five:eq:gramexception}
 4d_a C_{pq}C_{pr}C_{qr}.
\end{equation}
These are direct determinant expansions of $M$.

If another exponent $r$ lies strictly between $p$ and $q$, the nested
matching $(p,q),(a,r)$ is uniquely smallest by
\eqref{five:eq:onecollision}--\eqref{five:eq:convex4}, with valuation $d+w_{ar}$.
The additional term has valuation
$2\epsilon+2d+w_{pr}+w_{qr}>2d+2w_{ar}$ by \eqref{five:eq:midpoint}.
The square of that matching product is the unique minimum determinant
term, a contradiction.

Thus $p,a,q$ are consecutive. If the other exponents lie on opposite
sides, write $r<p<a<q<s$. On the four null vectors $r,p,q,s$ the nested
matching $(r,s),(p,q)$ is uniquely smallest, contradicting
\eqref{five:eq:gramnull}.

Otherwise reflect the exponents so that $p<a<q<r<s$.
The four-null identity on $p,q,r,s$ forces
\[
 d+w_{rs}=w_{ps}+w_{qr},
\]
since the middle matching is strictly larger than the nested matching.
On $a,q,r,s$ the nested matching $(a,s),(q,r)$ is uniquely smallest.
The additional determinant term must have exactly its doubled valuation:
\[
 2\epsilon+d+w_{qr}+w_{qs}+w_{rs}=2w_{as}+2w_{qr}.
\]
Substitution gives $2\epsilon+w_{ps}+w_{qs}=2w_{as}$, contradicting
\eqref{five:eq:midpoint}. This proves $V\leq8$.
\end{proof}

\subsection{Sharpness with eight unit edges at every prime}

Set $c_3=(0,2,1)$ and $c_i=t_i(-s_i^2,1,s_i)$ for $i=0,2,4,6$.
In the dyadic case take
\[
 (s_0,s_2,s_4,s_6)=(16,80,2,10),\qquad
 (t_0,t_2,t_6)=(64,1/2,1),
\]
and define
\[
 t_4=\frac{1-t_0t_6(s_0-s_6)^2-4}{t_2(s_2-s_4)^2}
     =-\frac{769}{1014}.
\]
For an odd prime put $T=p$ and take instead
\[
 (s_0,s_2,s_4,s_6)=(T^6,T^6+T^7,T,T+T^2),\qquad
 (t_0,t_2,t_6)=(T^5,T^{-2},T),
\]
with the same formula for $t_4$ replacing $4$ by $T^2$.
The numerator is a unit. The output coefficient at exponent $6$
is exactly $p^2$, and each other nonzero coefficient has one contribution.
For both recipes the complete height list, at exponents $2,\ldots,10$, is
\[
 17,11,7,4,2,1,1,2,5.
\]
The slopes $-6,-4,-3,-2,-1,0,1,3$ increase strictly.
Thus $V=U=8$. Apply the isometry \eqref{eq:isometry} below with
$a=s_0,b=s_2$, and clear a common denominator.
Two coefficients become axis vectors and the remaining three have
three coordinates, giving $S=11$ and one integer-square norm.

For $p=2$ this gives the explicit 26-bit integer example
\begin{align*}
 P&=519168x^2+121680x^3-37681x^4+9126x^6,\\
 Q&=-66453504-3204240x^3+1169649x^4-1242150x^6,\\
 R&=-624624x^3+209937x^4-106470x^6.
\end{align*}
The sole nonzero norm is $16224^2$, at exponent $3$.
Its output heights are the displayed list shifted upward by ten.
Every unit edge has a simple linear residual factor, so each example
has exactly eight simple nonzero $\Q_p$ roots.


For $n\le4$, there are at most $\binom n2+1$ contributions when $k=1$,
so $V\le\binom n2\le2n-2$. Together with the five-support bound and
the six-support construction in Section~\ref{sec:sharpness}, this shows
that six is the smallest vector-support size for which $V>2n-2$ is
possible over $\Q_p$, for every prime $p$.

\section{The finite necessary-condition certificate}

\subsection{Complete geometry and vertex patterns}

The exceptional diagonal exponent is already a pair exponent. If
$V\ge12$, at least thirteen distinct pair sums are needed, among the
fifteen pairs $i<j$. Equality of distinct pair sums can occur only as
$A_i+A_l=A_j+A_k$ with $i<j<k<l$. There are fifteen possible elementary
equalities. Every pair partition with at least thirteen groups is
generated by at most two such equalities: choose a spanning tree within
each group. Taking all equivalence closures gives 119 partitions.
There are ten reflected primary placements, hence 1,190 geometry cases.

The geometry certificate accounts for all cases as follows.
\begin{center}
\begin{tabular}{lrp{7.1cm}}
\toprule
Outcome&Cases&Exact evidence\\
\midrule
Forced unlisted collision&915&Rational row-space identity between two
pair sums declared distinct\\
Impossible strict order&48&Integer linear contradiction from the
support equalities and $A_{i+1}>A_i$\\
Feasible geometry&227&Rational ordered support with exactly the
specified pair-sum equalities\\
\bottomrule
\end{tabular}
\end{center}
In a feasible partition with $N$ output groups, mark thirteen strict
vertices and leave $N-13$ groups unmarked. The unmarked groups are
allowed to be vertices too. This yields 1,050 marked cases with no
pair collisions, 1,176 with one, and 133 with two: 2,359 in total.
The classification uses no numerical bound on the exponents.

\subsection{Convexity rules}

Let $\psi$ be the lower Newton boundary after normalization. Extend
it outside the output interval with affine rays of slopes strictly
below the first edge slope and strictly above the last. Choose the
left slope nonpositive and the right slope nonnegative. Then all
marked vertices, including endpoints, remain strict vertices of the
extension, and $\psi$ stays positive. Set
$h_{ij}=\psi(A_i+A_j)$ for $0\le i\le j\le5$.

At marked groups, $H_e=\psi(e)$; at unmarked groups, $H_e\ge\psi(e)$.
Identically equal exponents have equal $h$ values, and $h_{pq}>0$.
The following rules are immediate from convexity:
\begin{align}
 h_{ij}+h_{kl}&\ge h_{il}+h_{kj} &&(i<k,\ j<l),\label{eq:monge}\\
 h_{pr}+h_{qr}&\ge2h_{ar}.\label{eq:midpoint}
\end{align}
The first is strict if either inner argument is marked; the second
is strict if $A_a+A_r$ is marked. For the first, both inner arguments
are strictly between the outer arguments and their sums agree.

We also use all equal-barycenter relations implied by the primary and
collision equalities. If $u+v=s+t$ and $u<s,u<t$, convexity gives
\begin{equation}\label{eq:barycenter}
 \psi(u)+\psi(v)\ge\psi(s)+\psi(t),
\end{equation}
strictly if either inner argument is marked. This is encoded as the
finite disjunction
\[
 u\ge s\quad\text{or}\quad u\ge t\quad\text{or}\quad
 \psi(u)+\psi(v)\ge\psi(s)+\psi(t).
\]
The support relations are checked by exact rational row reduction.
The inequalities are not inferred from one sampled support realization.

\begin{lemma}[Rigidity of two additive height equations]\label{lem:rigidity}
Let $\phi$ be finite and convex on an interval containing $[0,a+b]$,
where $s>0$, $a>s$, and $b>2s$. If
\[
 \phi(a+s)-\phi(a)=\tfrac12\phi(2s),\qquad
 \phi(a+b)=\phi(a)+\phi(b),
\]
then $\phi(t)=\lambda t$ on $[\min(a,2s),a+b]$, where
$\lambda=\phi(2s)/(2s)$.
\end{lemma}
\begin{proof}
Subtract $\lambda t$ and write the resulting convex function as $\chi$.
Then $\chi(2s)=0$, $\chi(a+s)=\chi(a)$, and
$\chi(a+b)=\chi(a)+\chi(b)$.
If $a\le b$, monotonicity of equal-length increments gives
$\chi(b+s)\ge\chi(b)$, while convexity on the equal-barycenter pairs
$(2s,a+b)$ and $(a+s,b+s)$ gives the reverse inequality.
Equality forces affinity on $[2s,a+b]$. The equal values at $b,b+s$
force slope zero, hence value zero. The equality $\chi(a)=0$ extends
this interval to $a$ if $a<2s$.
If $a>b>2s$, comparison with the zero secant on $[a,a+s]$ gives
$\chi(a)\le0$, and convexity gives $\chi(b)\le0$. Since $b>s$,
convexity also gives $\chi(a+b)\ge\chi(a)$, hence $\chi(b)\ge0$.
Thus $\chi(b)=\chi(a)=0$. Equality on the chord with endpoints
$2s,a+b$ at the interior point $b$ forces affinity and value zero
throughout that interval.
\end{proof}

\subsection{Coefficient and Gram rules}

For a coefficient group not containing the primary, the minimum among
its pair-entry heights and its output height occurs at least twice.
In particular, a singleton group has $H=w$. The primary group already
has two zero-height contributions, its pair and the exceptional
diagonal; we impose no further coefficient rule there.

For four ordered indices $i<j<k<l$, put
$X=C_{ij}C_{kl}$, $Y=C_{ik}C_{jl}$, and $Z=C_{il}C_{jk}$.
Their Gram determinant is
\begin{equation}\label{eq:gram}
 X^2+Y^2+Z^2-2XY-2XZ-2YZ=0
\end{equation}
if all four vectors are null. If the exceptional vector is present,
add $4d_a C_{rs}C_{rt}C_{st}$, where $r,s,t$ are the other indices.
Thus the minimum among
\begin{equation}\label{eq:minimum}
 2(w_{ij}+w_{kl}),\quad2(w_{ik}+w_{jl}),\quad2(w_{il}+w_{jk})
\end{equation}
and, in the exceptional case,
\[
 2\epsilon+w_{rs}+w_{rt}+w_{st}
\]
must occur at least twice. A unique minimum in this list would remain
unique in the full determinant, because a mixed term has height
$\epsilon$ plus the average of two square-term heights and
$\epsilon\ge0$. The argument remains valid in residue characteristic two.

For an exceptional quartet whose six pairs are singleton marked groups
and avoid the primary, strict convexity makes the nested matching
uniquely smallest among the three matchings. The additional diagonal
term must tie its square exactly. This forced linear equality is the
certificate rule named \texttt{clean\_gram}.

\subsection{Direct obstructions and finite disjunction trees}

Let $L$ contain the primary and all pairs in nonsingleton groups, and
let $U_0$ contain all pairs in nonsingleton or unmarked groups. A null
quartet $i<j<k<l$ is impossible whenever
\[
 \{ij,kl,ik,jl\}\cap L=\varnothing,
 \qquad \{il,jk\}\cap U_0=\varnothing.
\]
Indeed, the nonnested entries are at least their boundary heights,
while the nested entries are at most their boundary heights: they
equal the marked singleton heights, or are the primary at zero below
its positive boundary. The nested arguments are marked strict
vertices. Strict convexity therefore makes the nested matching
uniquely smallest, contrary to \eqref{eq:gram}. This excludes 1,397
marked cases.

Each of the remaining 962 cases is encoded by the preceding linear
rules and finite disjunctions. A repeated minimum is encoded by choosing
every possible pair of equal minimal entries. At an internal proof-tree
node, every branch of the selected disjunction is present. A leaf gives
integer multipliers for its accumulated linear constraints. Multipliers
of inequalities are nonnegative, multipliers of equalities are unrestricted,
and some strict inequality has positive weight. The weighted sum of
the linear forms is identically zero, giving $0>0$.

The certificate has 9,085 nodes and 6,627 such leaves. Exactly six cases
require the additional equal-barycenter rules. The solver-free checker
independently enumerates the partitions, verifies all geometry certificates,
regenerates the marked patterns, checks all direct null obstructions,
audits each rule instance, verifies every branch, and evaluates each
integer identity exactly. All 2,359 cases are excluded. This proves
the all-edge assertion of Theorem~\ref{thm:main}.

\section{The unit-edge bound}

Suppose $U\ge11$. The established $V\le11$ gives $V=U=11$, so the
twelve strict vertices have consecutive integer exponents and the
output span is eleven. Apply Lemma~\ref{lem:perturb}. All old vertices
persist, and the bound $V\le11$ permits no additional vertex. The
output interval is therefore unchanged. Since all pair entries are
now finite and the hidden diagonal exponent is a pair sum,
\[
 A_4+A_5-A_0-A_1=11.
\]
Translate $A_0$ to zero. The integer support order gives
$A_4-A_1\ge3$, whence $A_5\le8$. This bound follows from the
assumed eleven unit edges; it is not an imposed search limit.

Exactly four six-element subsets of $\{0,\ldots,8\}$ containing zero
have twelve consecutive pair sums:
\[
 \begin{split}
 &\{0,1,2,3,4,8\},\qquad\{0,2,3,4,5,8\},\\
 &\{0,3,4,5,6,8\},\qquad\{0,4,5,6,7,8\}.
 \end{split}
\]
After reflection $a\in\{1,2\}$, and there are ten possible primary
placements across these supports. Mark every one of the twelve output
groups. The same linear rules exclude all ten cases, with 749 nodes
and 567 exact contradiction leaves; two cases use equal-barycenter
convexity. The separate checker \texttt{verify\_unit\_six.py} regenerates
the supports and primary placements and verifies the full certificate.
This proves $U\le10$.

\section{Simultaneous sharpness and a compact formula}\label{sec:sharpness}

Set $c_3=(0,2,1)$ and $c_i=t_i(-s_i^2,1,s_i)$ for
$i\in\{0,2,4,6,8\}$. Over $\Q_p$ put $T=p$ and use the parameters
in Table~\ref{tab:parameters}.
\begin{table}[ht]
\centering
\begin{tabular}{crrrrrrrrr}
\toprule
Case&$A$&$B$&$D$&$E$&$u$&$v$&$w$&$H_6$&$H_{10}$\\
\midrule
$p=2$&39&41&18&13&10&$-17$&$-2$&20&18\\
$p$ odd&33&34&14&11&9&$-14$&$-1$&16&14\\
\bottomrule
\end{tabular}
\caption{Parameters for the sharp six-support formula.}\label{tab:parameters}
\end{table}
Define
\begin{gather*}
 s_0=T^A,\quad s_2=T^A+T^B,\quad s_4=T^{-1},\quad
 s_6=T^D,\quad s_8=T^E,\\
 t_0=T^u,\quad t_2=T^v,\quad t_6=T^w,\\
 t_4=\frac{1-T^{H_6}-t_0t_6(s_0-s_6)^2}{t_2(s_2-s_4)^2},\qquad
 t_8=\frac{-T^{H_{10}}-t_4t_6(s_4-s_6)^2}{t_2(s_2-s_8)^2}.
\end{gather*}
The directions are distinct and all denominators are nonzero.
The coefficient norms are zero except $\qf(c_3)=1$.

The definitions force the output coefficients at 6 and 10 to be
$T^{H_6}$ and $T^{H_{10}}$. The numerator defining $t_4$ is a unit:
$H_6>0$ and $u+w+2D>0$. Hence $\val_p(t_4)=2-v$.
The two terms in the numerator defining $t_8$ have unequal valuations,
with $w-v<H_{10}$. Therefore
$\val_p(t_8)=w-2v-2E$.

With $\epsilon=\val_p(2)$, the output heights in exponent order
$2,3,\ldots,12,14$ are
\begin{gather*}
 u+v+2B,\quad \epsilon+u+A,\quad u-v,\quad \epsilon+v+A,\quad H_6,\\
 \epsilon-v,\quad v+w+2D,\quad \epsilon+w+D,\quad H_{10},\\
 \epsilon+w-2v-E,\quad w-3v-2E,\quad 2w-2v.
\end{gather*}
Every unforced coefficient has a unique minimum contribution. At
exponent 8 the two heights are $u+w-2v$ and $v+w+2D$, and the latter
is smaller because $u-3v>2D$. All other unforced coefficients have a
single pair. The resulting tables are
\begin{center}\small
\begin{tabular}{crrrrrrrrrrrr}
\toprule
Exponent&2&3&4&5&6&7&8&9&10&11&12&14\\
\midrule
$p=2$&75&50&27&23&20&18&17&17&18&20&23&30\\
$p$ odd&63&42&23&19&16&14&13&13&14&16&19&26\\
\bottomrule
\end{tabular}
\end{center}
For odd $p$ the consecutive slopes are
\[
 -21,-19,-4,-3,-2,-1,0,1,2,3,7/2;
\]
for $p=2$ the first two become $-25,-23$. They are strictly increasing.
Thus all twelve points are strict vertices, the first ten edges have
width one, and the last has width two. This proves $(V,U)=(11,10)$.

For integer coordinates, put $a=s_0$, $b=s_2$, and apply
\begin{equation}\label{eq:isometry}
 (P,Q,R)\longmapsto
 \frac{(P+2aR-a^2Q,\ -P-2bR+b^2Q,\ -P-(a+b)R+abQ)}{a-b}.
\end{equation}
Direct expansion verifies preservation of $PQ+R^2$. The vectors at
0 and 2 become coordinate-axis vectors; the other four have three
nonzero coordinates for the displayed rational parameters. Clearing a
common denominator $M$ gives $S=1+1+4\cdot3=14$, a sole nonzero norm
$M^2$, and a common output-height shift $2\val_p(M)$. The dyadic
frozen witness has maximum input length 228 bits.

If $\val(2)=0$ in a general nontrivially valued field, choose any $T$
with positive valuation and use the odd row of the table. The displayed
heights are multiplied by $\val(T)$. The valuation proof uses only
unit leading coefficients. Independently, the exact formal checker
clears denominator factors with constant term one and expands in
$\Z[T,T^{-1}]$; every leading coefficient is a signed power of two.
This verifies specialization in every odd positive characteristic too.

If $\val(2)>0$, then $K$ has characteristic zero and $\val|_{\Q}$ is
$\val(2)$ times the ordinary dyadic valuation. Indeed, every odd
integer is a unit by the ultrametric inequality. The rational dyadic
formula therefore specializes to $K$. This proves sharpness in
Theorem~\ref{thm:main}. A trivial valuation has at most one strict
lower edge, explaining the nontriviality qualification.

Finally, over $\Q_p$ every unit edge has a linear residual polynomial
with a nonzero simple root. Hensel's lemma gives one simple root of
its valuation. The last edge has nonintegral slope $7/2$ and gives no
$\Q_p$ root. This proves Corollary~\ref{cor:prime}.

\section{An all-size one-norm cancellation family}\label{sec:family}

We prove the following uniform construction directly.
\begin{theorem}\label{thm:family}
For every prime $p$ and integer $n\ge8$, there are integer polynomials
$P,Q,R$ with
\[
 k=1,\qquad V_p=2n+1,\qquad U_p=2n,\qquad S=3n-4.
\]
Their vector support is $\{0,1,2,4\}\cup\{7,\ldots,n+2\}$.
For each $p$, one fixed nonzero integer-square norm works for every $n$.
The input coefficient bit lengths are $O(n^2\log p)$, and the output
has exactly $2n$ simple nonzero $\Q_p$ roots, of distinct valuations.
\end{theorem}

\subsection{The eight-support core}
Set $T=p$, $c_4=(0,2,1)$, and $c_i=t_i(-s_i^2,1,s_i)$ elsewhere.
In both residue characteristics put
\[
 s_0=T^{-3},\qquad s_1=T^{-4},\qquad s_2=1+T^{35},
 \qquad s_9=1+T^{152},\qquad s_{10}=1+T^{144}.
\]
For $p=2$ take
\[
\begin{split}
 s_7&=1+T^{141},\qquad s_8=1+(T^{18}-1)T^{142},\\
 t_0&=T^{226},\qquad t_1=\frac{T^{153}}{(1-T^4)^2},
 \qquad t_2=T^{74}.
\end{split}
\]
For odd $p$ take
\[
\begin{split}
 s_7&=1+(T-1)T^{141}/2,\qquad
 s_8=1+(T^{18}-1)T^{142}/2,\\
 t_0&=T^{225},\qquad t_1=\frac{T^{152}}{(1-T^4)^2},
 \qquad t_2=T^{73}.
\end{split}
\]
Write $D_{ij}=(s_i-s_j)^2$ and define $t_7,t_8,t_9,t_{10}$ by
\begin{equation}\label{eq:family-system}
\begin{pmatrix}
t_1D_{17}&t_0D_{08}&0&0\\
t_2D_{27}&t_1D_{18}&t_0D_{09}&0\\
0&t_2D_{28}&t_1D_{19}&t_0D_{0,10}\\
-2s_7(1-s_7)&0&t_2D_{29}&t_1D_{1,10}
\end{pmatrix}
\begin{pmatrix}t_7\\t_8\\t_9\\t_{10}\end{pmatrix}
=\begin{pmatrix}1-T^{61}\\-T^{48}\\-T^{36}\\-T^{25}\end{pmatrix}.
\end{equation}
The null-null cross coefficients are $-t_it_jD_{ij}$, and those with
$c_4$ are $2t_is_i(1-s_i)$. Thus these equations prescribe the output
coefficients at $8,9,10,11$ exactly.
The dyadic valuation matrix is
\[
\begin{pmatrix}
145&220&\infty&\infty\\
144&145&220&\infty\\
\infty&144&145&220\\
142&\infty&144&145
\end{pmatrix};
\]
for odd $p$, subtract one from every finite entry.
The diagonal determinant term is uniquely minimal, of valuation
$580$ or $576$: every nontrivial permutation cycle uses an upper
entry with excess $75$, whereas the total decreases from the other
entries in that cycle are at most five. Put $K=145$ or $144$, respectively, and normalize
\[
 a=T^Kt_7,\quad b=T^{K+1}t_8,\quad
 c=T^{K+2}t_9,\quad d=T^{K+2}t_{10}.
\]
Multiplying the rows by $1,T,T^2,T^3$, respectively, the first
three normalized rows give
\[
 a=1-T^{61}\pmod {T^{74}},\quad
 b=-a-T^{49}\pmod {T^{74}},\quad
 c=-b-T^{38}\pmod {T^{75}}.
\]
The normalized $t_1D_{1j}$ correction is
$((1-T^4s_j)/(1-T^4))^2$, which is one to precision greater than
$140$. The normalized $D_{2j}$ correction is one to precision at
least $106$. These exceed all precisions used in the elimination.
The fourth row gives
\[
 a+c+2d=-2^{28}\pmod {2^{49}}\quad(p=2),
 \qquad
 (T-1)a+c+Td=-T^{28}\pmod {T^{49}}\quad(p\ne2).
\]
The normalized matrix is integral with determinant valuation one.
Its cofactor formula first bounds all four unknown valuations below by
$-1$. The first three rows then give their congruences with precision
reduced by one; the fourth gives $Td=-T\pmod {T^{28}}$. This proves
integrality and recovers the full displayed precisions.
Thus $a,c\equiv1$ and $b,d\equiv-1\pmod {T^{27}}$.
The coefficient at exponent $12$ is
\[
 -T^{-3}\left((T^{18}-1)s_8b+
       \left(\frac{s_2-s_{10}}{T^{35}}\right)^2d\right),
\]
whose bracket is $-T^{18}$ modulo $T^{27}$.
The four solved coefficients and all remaining unique minima now give
the strict dyadic vertices
\[
\begin{split}
 &(1,371),(2,294),(3,219),(5,146),(6,110),(7,75),\\
 &(8,61),(9,48),(10,36),(11,25),(12,15),(13,6),(14,-2),\\
 &(15,-9),(16,-10),(17,-10),(18,-9),(19,-6).
\end{split}
\]
For odd $p$, replace these boundary heights as follows:
\[
\begin{array}{c|rrrrr}
e&1&2&3&5&6\\ \hline
\val(f_e)&369&292&217&144&108\\
e&15&16&17&18&19\\ \hline
\val(f_e)&-7&-8&-8&-7&-4
\end{array}
\]
The sole additional nonzero point has exponent $4$ and height
$221$ or $219$, respectively, strictly above the boundary.
The dyadic slopes are
\[
\begin{split}
&-77,-75,-73/2,-36,-35,-14,-13,-12,-11,\\
&-10,-9,-8,-7,-1,0,1,3.
\end{split}
\]
At odd primes, $-35$ becomes $-33$, and $-7$ becomes $-5$.
In either case there are seventeen strict edges and sixteen unit
edges. The sole nonunit edge, from exponent $3$ to $5$, has slope
$-73/2$.

\subsection{The tail and its complete height comparison}

For $11\leq j\leq n+2$ append the dyadic tail
\[
 t_j=T^{4j^2-66j+113},\qquad
 s_j=1+T^{144}+\sum_{r=11}^jT^{6r+89},
\]
or the odd-prime tail
\[
 t_j=T^{4j^2-66j+114},\qquad
 s_j=1+T^{144}+\sum_{r=11}^jT^{6r+88}.
\]
Put $h_j=\val_p(t_j)$, using the displayed quadratic also at $j=10$.
For $10\le i<j$, the direction distance has valuation $6i+95$ in the
dyadic case and $6i+94$ otherwise. In both cases the cross height is
\[
 4i^2+4j^2-54i-66j+416.
\]
At a fixed sum $e=i+j\ge21$, this is a strictly convex quadratic
in $i$, with real minimizer $e/2-3/4$. The unique allowed integer
minimizer is the adjacent pair for odd $e$ and the gap-two pair for
even $e$. Its height is
\begin{equation}\label{eq:family-tail-height}
 H_e=2e^2-60e+412.
\end{equation}
The pair $(9,11)$ has height $12$ for $p=2$ and $14$ otherwise, at
exponent $20$. More generally, the height of $(9,j)$ minus $H_{j+9}$
is $2(j-10)(j-11)$ for $p=2$, and that expression plus two otherwise.
It is positive for every $j\ge12$.

For the other old indices, the new cross heights are $h_j+C_i$, where
\[
\begin{array}{c|rrrrrr}
i&0&1&2&4&7&8\\ \hline
C_i\ (p=2)&220&145&144&145&137&138\\
C_i\ (p\ne2)&219&144&143&144&138&139
\end{array}
\]
For $p=2$, subtracting $H_{j+i}$ gives
\[
\begin{array}{c|l|c}
i&\text{difference}&\text{range for }j+i\ge20\\ \hline
0&2j^2-6j-79&j\ge20\\
1&2j^2-10j-96&j\ge19\\
2&2j^2-14j-43&j\ge18\\
4&2j^2-22j+54&j\ge16\\
7&2j^2-34j+160&j\ge13\\
8&2j^2-38j+191&j\ge12
\end{array}
\]
For odd $p$, the first four differences are unchanged and the last
two increase by two. Each polynomial is positive and increasing in
the indicated range. For the finitely many $j\ge11$ with $j+i<20$,
substitution in the core table also gives strict inequalities; the
smallest margin is one, at $(i,j)=(1,11)$. At exponent $20$, every
competitor to $(9,11)$ is strictly higher, including in the odd case
where the boundary height is $14$ instead of $H_{20}=12$.

Hence no core vertex is lost and all new coefficients have a unique
minimum contribution. The added vertices are $(20,12)$ for $p=2$,
$(20,14)$ for odd $p$, and $(e,H_e)$ for $21\le e\le2n+3$.
The first new slopes are $18,22,26,30,\ldots$ for $p=2$ and
$18,20,26,30,\ldots$ for odd $p$. They increase strictly, starting
above the core's final slope three. Each appended coefficient therefore
adds two unit edges. This proves $V=17+2(n-8)=2n+1$ and
$U=16+2(n-8)=2n$ for all $n\ge8$.

\subsection{Integer coefficients, sparsity, and roots}

For the integer realization, put $a=s_0,b=s_1$ and apply
\begin{equation}\label{eq:strongeriso}
(p_0,q_0,r_0)\longmapsto
\frac{(p_0+2ar_0-a^2q_0,\,
       -p_0-2br_0+b^2q_0,\,
       -p_0-(a+b)r_0+abq_0)}{a-b}.
\end{equation}
Direct expansion proves that this is a $\qf$-isometry.
The first two coefficient vectors become coordinate axes, and all
others have three nonzero coordinates. Hence $S=3n-4$.
Each matrix denominator divides $p^4(p-1)$ and every new scale
denominator divides $p^{129}$. The core denominator already contains
$p^{K+6}$, from its exponent-$9$ coefficient, and contains $p-1$,
from the exponent-$2$ coordinate
$p^{\val(t_2)-4}(p^4+p^{39}-1)^2/(p-1)$.
Its numerator is coprime to $p-1$.
Thus the same core denominator clears every tail, giving one fixed
integer-square norm. The displayed exponents give $O(n^2\log p)$ input bits.
This completes Theorem~\ref{thm:family}.

Each unit edge gives one simple nonzero $\Q_p$ root, by its linear
residual polynomial. The remaining slope $-73/2$ is nonintegral
at every prime. Consequently the nonzero $\Q_p$ root count is
exactly $2n$.



In particular, the candidate inequalities $V\le2n+2k-4$ and
$U\le2n+2k-4$ fail already at $k=1$. The exceptional coefficient in
this construction is interior, at exponent four. The family does not
represent a prescribed target polynomial with an exceptional constant term.

\section{Computational verification}

The ancillary archive \texttt{verification.zip} contains the exact
certificates, example coefficients, and checking programs. After
extracting the archive, \texttt{python -B -X utf8 verify.py} checks
the file manifest and runs all five verification programs. The
programs can also be run individually:
\begin{verbatim}
python -B -X utf8 verify_sharp_six.py
python -B -X utf8 verify_unit_six.py
python -B -X utf8 verify_unit10_examples.py
python -B -X utf8 verify_positive_control.py
python -B -X utf8 verify_families.py
\end{verbatim}
They require no solver. The first two regenerate the rule systems,
independently audit the formula and justification metadata of every
instance, and check all finite coverage and integer contradictions.
The third imports no constructor: it expands the frozen integer
inputs by ordered scalar convolution, checks every coefficient
against every proposed supporting line, checks all coefficient norms,
and lifts sixty roots modulo $p^{24}$ at
$p=2,3,5,7,11,17$. It also verifies the formal field-uniform sharpness
calculation. The fourth evaluates 192 axioms and 4,176 disjunctions
on a genuine eleven-edge equality example and rejects invalid
arithmetic-certificate controls.

The fifth checker verifies the results in
Sections~\ref{sec:five} and~\ref{sec:family}: six frozen five-support
examples, twenty-five family members through $n=35$ at six primes,
and 144 simple root lifts to precision $p^{24}$. It checks the twelve
polynomial tail inequalities on their full integer ranges using
their coefficients and monotonicity, the sixty-four finite core
comparisons, and uniqueness of the core determinant minimum in both
residue-characteristic cases. The general upper bound at five supports
and the uniform family assertions follow from the written proofs.

Z3 was used to generate the certificates. The checking programs
verify their integer identities directly and require only the Python
standard library. Each selected proof file is hashed in a manifest;
the checkers validate these hashes and record their own source hashes.
Python optimization flags are rejected because the implementation
uses assertions for exact checks. The archive is self-contained.
The mathematical reduction to the verified finite systems is proved
in the preceding sections.

\section{Further questions}

The convexity argument uses additive relations among exponent
sums before introducing their boundary heights. It resolves six
patterns that the principal Gram conditions and the basic Monge and
primary-midpoint inequalities did not exclude. The resulting finite
proof depends on at most two pair-sum losses when thirteen vertices
are demanded from fifteen pairs. This counting feature is specific
to the present extremal problem. No extrapolation to arbitrary support
size follows from the certificate.

The next broader questions concern a uniform bound for growing $n$
with one nonzero norm, and restrictions on endpoint-normalized
representations of specified completely convex outputs. Theorem~\ref{thm:family} rules out the candidate bound $2n+2k-4$
for both edge statistics. The sharp five- and six-support results
and this all-size construction leave open a uniform upper bound
for growing support. None of these results supplies an unrestricted
circuit separation.

\paragraph{Acknowledgment.}
The author thanks OpenAI\textquoteright{}s GPT-6 for assistance with
exploratory calculations, proof development, code checks, and manuscript
drafting. The author takes responsibility for the mathematical claims
and the final manuscript.

\begin{thebibliography}{1}
\bibitem{KPTT}
P. Koiran, N. Portier, S. Tavenas, and S. Thomass\'e,
\emph{A tau-conjecture for Newton polygons}, arXiv:1308.2286v2 (2014).
\url{https://arxiv.org/abs/1308.2286v2}

\end{thebibliography}
\end{document}
